\subsection{平方可积表示}

在本目中，所考虑的希尔伯特空间都是复的。假设 G 是单模的，其中心记作 C。对 C 的每个闭子群 Z，用 \(\beta_{Z}\) 表示 Z 上的一个哈尔测度，并给 G/Z 赋以哈尔测度 \(\nu_{Z} = \mu/\beta_{Z}\)（INT, VII, 第 44 页, § 2, 第 2 目, 定义 1）。

设 \(\pi\) 是 G 在希尔伯特空间 E 中的一个不可约酉表示，\(\chi \in \widehat{\mathbf{C}}\) 是其中心特征标（V, 第 390 页, 定义 6）。对 E 中所有 \(x\) 与 \(y\)，用 \(f_{x,y}\) 表示矩阵系数 \(g \mapsto \langle x \mid \pi(g)y \rangle\)；它是 G 上取复值的连续函数。

设 E 中的 \(x\) 与 \(y\)。对 \(z \in \mathrm{C}\) 与 \(g \in \mathrm{G}\)，我们有

\[
f _ {x, y} (z g) = \langle x \mid \pi (z g) y \rangle = \langle x \mid \chi (z) \pi (g) y \rangle = \chi (z) f _ {x, y} (g),\tag{9}
\]

因此 \(f_{x,y}\) 属于空间 \(\mathcal{F}_{\chi}(\mathrm{G})\)（V, 第 408 页）。此外，对每对 \((g_1, h_1) \in \mathrm{G} \times \mathrm{G}\) 与每个 \(g \in \mathrm{G}\)，我们有

\[
f _ {\pi (g _ {1}) x, \pi (h _ {1}) y} (g) = \langle \pi (g _ {1}) x | \pi (g) \pi (h _ {1}) y \rangle = f _ {x, y} (g _ {1} ^ {- 1} g h _ {1}).\tag{10}
\]

关系式 (9) 为下面的定义提供了依据。

定义 3. --- 设 \(\pi\) 是 G 在希尔伯特空间 E 中的一个不可约酉表示。我们称 \(\pi\) 是模中心平方可积的，如果由函数 \(|f_{x,y}|\) 通过取商得到的 G/C 上的函数属于 \(\mathcal{L}^2 (\mathrm{G / C})\)，对所有 \((x,y)\in \mathrm{E}\times \mathrm{E}\) 成立。

这一条件不依赖于 G/C 上哈尔测度的选取。

若 \(\pi\) 的矩阵系数都属于 \(\mathcal{L}^2 (\mathrm{G})\)，则称 \(\pi\) 是平方可积的；G 的一个平方可积不可约酉表示的存在蕴含 G 的中心是紧的（V, 第 487 页, 习题 5）。

存在这样的群 G，它没有任何平方可积表示，即使模中心也没有（参见 V, 第 516 页, 习题 32）。

命题 6. --- 设 \(\pi\) 是 G 在希尔伯特空间 E 中的一个不可约酉表示。设 \(\chi\) 是 \(\pi\) 的中心特征标。则 \(\pi\) 模中心平方可积当且仅当存在 E 中非零元素 \(x_0\) 与 \(y_0\)，使得由 \(|f_{x_0,y_0}|\) 通过取商诱导的 G/C 上的函数属于 \(\mathcal{L}^2 (\mathrm{G / C})\)。

设存在 E 的非零元素 \(x_{0}\) 与 \(y_{0}\)，使得 \(|f_{x_{0},y_{0}}| \in \mathcal{L}^{2}(\mathrm{G/C})\)。只需证明这时 \(\pi\) 模中心平方可积。

 设 F 是由元素 \(x \in E\)——这里 \(|f_{x,y_{0}}|\) 属于 \(\mathcal{L}^{2}(G/C)\)——组成之集。这是 E 的一个向量子空间；它包含 \(x_{0}\)，从而非零。上面的关系式 (10) 蕴含 F 在 \(\pi\) 下稳定；由于表示 \(\pi\) 不可约，空间 F 在 E 中稠密。

 设 \(x \in \mathrm{F}\)。由于 \(f_{x,y_0}\) 属于 \(\mathcal{F}_{\chi}(\mathrm{G})\) 且 \(\mu\) -可测，又 \(|f_{x,y_0}| \in \mathcal{L}^2(\mathrm{G/C})\)，我们有 \(f_{x,y_0} \in \mathcal{L}_{\chi}^2(\mathrm{G})\)（V, 第 413 页的命题 5 应用于 Z = C）。用 \(u\) 表示 E 到 \(\mathrm{L}_{\chi}^{2}(\mathrm{G})\) 的部分算子，其定义域为 F，且把 \(x \in \mathrm{F}\) 映为 \(f_{x,y_0}\) 在 \(\mathrm{L}_{\chi}^{2}(\mathrm{G})\) 中的类。

 我们来证明部分算子 \(u\) 是闭的。设 \((x_{n}, u(x_{n}))_{n \in \mathbf{N}}\) 是 \(u\) 的图中的元素组成的序列，它在 \(\mathrm{E} \times \mathrm{L}_{\chi}^{2}(\mathrm{G})\) 中收敛。设 \(x\) 是 \((x_{n})\) 的极限，并设 \(f \in \mathcal{L}_{\chi}^{2}(\mathrm{G})\) 是一个函数，其类是序列 \((u(x_{n}))\) 的极限。

函数 \(u(x_{n})\) 是矩阵系数 \(f_{x_{n},y_{0}}\) 的类。对每个 \(g \in G\)，我们有

\[
f _ {x _ {n}, y _ {0}} (g) = \langle x _ {n} | \pi (g) y _ {0} \rangle \rightarrow \langle x | \pi (g) y _ {0} \rangle = f _ {x, y _ {0}} (g)
\]

当 \(n \to +\infty\)。因此我们有 \(f_{x,y_{0}} \in \mathcal{L}_{\chi}^{2}(\mathrm{G})\)，且模 C 几乎处处有 \(f = f_{x_{0},y}\)（V, 第 409 页, 命题 3 的推论）；这就是说 \(u(x)\) 是 f 在 \(\mathrm{L}_{\chi}^{2}(\mathrm{G})\) 中的类。于是我们证明了 u 是闭的。

u 的定义域 F 在 \(\pi\) 下稳定，且关系式 (10) 表明 u 满足 \(u \circ \pi(g) = \gamma_{\mathrm{G},\chi}(g) \circ u\)，对所有 \(g \in G\)。从而我们也有等式 \(u^{*} \circ \gamma_{\mathrm{G},\chi}(g) = \pi(g) \circ u^{*}\)，对所有 \(g \in G\)（V, 第 381 页, 引理 6），由此得 \((u^{*} \circ u) \circ \pi(g) = \pi(g) \circ (u^{*} \circ u)\)。而部分算子 \(u^{*} \circ u\) 是自伴的（IV, 第 241 页, 命题 12），特别地是闭的（IV, 第 236 页, 命题 7），于是 V, 第 387 页, 推论 1 蕴含 \(u^{*} \circ u\) 的定义域等于 E。更有，我们有 F = E，也就是说，函数 \(|f_{x,y_{0}}|\) 属于 \(\mathcal{L}^{2}(\mathrm{G}/\mathrm{C})\)，对所有 \(x \in E\) 成立。

设 \((x,y)\in\mathrm{E}\times\mathrm{E}\)。我们有 \(f_{x,y_{0}}\in\mathcal{L}_{\chi}^{2}(\mathrm{G})\)。把 \(\delta_{G,\chi}\) 代替 \(\gamma_{G,\chi}\)，用同样方法可证：\(y\in E\) 中使函数 \(|f_{x,y}|\) 属于 \(\mathcal{L}^{2}(\mathrm{G}/\mathrm{C})\) 的那种元素所成之集等于 E。命题得证。

在本节余下部分，我们取定 C 的一个使 C/Z 为紧的闭子群 Z。

引理 9. --- 设 \(\pi\) 是 G 在希尔伯特空间 E 中的一个模中心平方可积的不可约酉表示。设 \(\chi\) 是 \(\pi\) 的中心特征标到 Z 的限制。对 E 中所有 \(x\) 与 \(y\)，我们有 \(f_{x,y} \in \mathcal{L}_{\chi}^{2}(\mathrm{G})\)。

 函数 \(f_{x,y}\) 连续，因此 \(\mu\) -可测。由公式 (9)（第 420 页）我们有 \(f_{x,y} \in \mathcal{F}_{\chi}(\mathrm{G})\)。此外，由 INT, VII, 第 64 页, § 2, 第 8 目, 推论 1, c)，我们有 \(\mathrm{N}_{2}(f_{x,y}) < +\infty\)，因为 C/Z 紧且 \(|f_{x,y}| \in \mathcal{L}^{2}(\mathrm{G}/\mathrm{C})\)。于是断言由 V 第 413 页的命题 5 得出。

命题 7. --- 设 π 是 G 在希尔伯特空间 E 中的一个模中心平方可积的不可约酉表示。设 χ 是 π 的中心特征标到 Z 的限制。

存在一个实数 c \textgreater{} 0 与一个唯一的 \((\mathbf{G} \times \mathbf{G})\) -等距态射 w，它把酉表示 \(\overline{\pi} \boxtimes \pi\) 映到 \(\mathrm{L}_{\chi}^{2}(\mathrm{G})\) 中，使得对所有 \((x, y) \in \overline{\mathrm{E}} \times \mathrm{E}\)，元素 \(w(x \otimes y)\) 是 \(\mathrm{L}_{\chi}^{2}(\mathrm{G})\) 中函数 \(c^{1/2} f_{x,y}\) 的类。

 对每对 \((x,y)\in\mathrm{E}\times\mathrm{E}\)，我们有 \(f_{x,y}\in\mathcal{L}_{\chi}^{2}(\mathrm{G})\)（引理 9）。用 v 表示 \(\overline{E}\otimes E\) 到 \(\mathrm{L}_{\chi}^{2}(\mathrm{G})\) 中的唯一线性映射，使得 \(v(x\otimes y)\) 是 \(f_{x,y}\) 的类，对所有 \((x,y)\in\overline{\mathrm{E}}\times\mathrm{E}\) 成立。

我们在下面证明如下的引理。

引理 10. --- 存在一个实数 c \textgreater{} 0，使得线性映射 \(w = c^{1/2}v\) 是等距的。

假设该引理成立，注意公式 (10)（第 420 页）可以写成

\[
v (\overline {{\pi}} (g _ {1}) x \otimes \pi (h _ {1}) y) = \varrho_ {\mathrm{G}, \chi} (g _ {1}, h _ {1}) v (x \otimes y)\tag{11}
\]

 对所有 \((g_{1}, h_{1}) \in \mathbf{G} \times \mathbf{G}\) 与每对 \((x, y) \in \overline{\mathbf{E}} \times \mathbf{E}\) 成立。\(\overline{\mathbf{E}} \otimes \mathbf{E}\) 到 \(\mathrm{L}_{\chi}^{2}(\mathrm{G})\) 中的等距线性映射 w 可连续延拓（仍记作 w）到 \(\overline{E} \widehat{\otimes}_{2} E\) 上。由连续性与线性性，公式 (11) 蕴含 w 是一个 \((\mathrm{G} \times \mathrm{G})\) -态射——它从 \(\overline{\pi} \boxtimes \pi\) 映到 \(\mathrm{L}_{\chi}^{2}(\mathrm{G})\) 中——这就完成了命题的证明。

 我们来证明该引理。对所有 \(x \in E\)，用 \(u_{x}\) 表示线性映射 \(y \mapsto v(x \otimes y) = f_{x,y}\)（它从 E 映到 \(\mathrm{L}_{\chi}^{2}(\mathrm{G})\) 中）。我们有 \(u_{x} \in \mathrm{Hom}_{\mathrm{G}}(\pi, \gamma_{\mathrm{G},\chi})\)。由公式 (10)，根据 V 第 388 页的推论 5，存在一个实数 \(\lambda_{x} \geqslant 0\)，使得 \(\lambda_{x}u_{x}\) 是等距的。

设 E 中的 \(x\) 与 \(y\)。我们有

\[
\| f _ {x, y} \| ^ {2} = \int_ {\mathrm{G/Z}} \overline {{f}} _ {x, y} f _ {x, y} d \nu_ {\mathrm{Z}} = \int_ {\mathrm{G/Z}} \check {\overline {{f}}} _ {x, y} \check {f} _ {x, y} d \nu_ {\mathrm{Z}}
\]

因为 G/Z 是单模的（V 第 417 页, 引理 7）。注意到 \(\check{f}_{x,y} = \overline{f}_{y,x}\)，我们得到

\[
\lambda_ {x} \| y \| ^ {2} = \| f _ {x, y} \| ^ {2} = \int_ {\mathrm{G/Z}} f _ {y, x} \overline {{f}} _ {y, x} d \nu_ {\mathrm{Z}} = \| f _ {y, x} \| ^ {2} = \lambda_ {y} \| x \| ^ {2}.
\]

这表明正实数 \(\lambda_{x}/\|x\|^{2}\) 与 E 中非零元素 x 的选取无关。它是严格正的，因为对每个非零的 x，函数 \(f_{x,x}\) 连续且在 e 处取值 \(\|x\|^{2}>0\)，由此 \(\|u_{x}(x)\|=\|f_{x,x}\|>0\)。用 \(c^{-1}\) 记这个实数。

对每对 \((x,y)\in\mathrm{E}\times\mathrm{E}\)，由此得出

\[
\| v (x \otimes y) \| ^ {2} = \| f _ {x, y} \| ^ {2} = \lambda_ {x} \| y \| ^ {2} = c ^ {- 1} \| x \| ^ {2} \| y \| ^ {2} = c ^ {- 1} \| x \otimes y \| ^ {2}.
\]

利用 EVT, V, 第 29 页, 推论 1，我们推出线性映射 \(w = c^{1/2}v\)（它从 \(\overline{\mathrm{E}}\widehat{\otimes}_2\mathrm{E}\) 映到 \(\mathrm{L}_{\chi}^{2}(\mathrm{G})\) 中）是等距的，这正是所要证的。

 推论. --- 设 \(\pi\) 是 G 的一个不可约酉表示，\(\chi\) 是其中心特征标到 Z 的限制。表示 \(\pi\) 模中心平方可积当且仅当它同构于表示 \(\gamma_{\mathrm{G},\overline{\chi}}\)（从 G 到 \(\mathrm{L}_{\overline{\chi}}^{2}(\mathrm{G})\) 中）的一个子表示（相应地，同构于表示 \(\delta_{\mathrm{G},\chi}\)（从 G 到 \(\mathrm{L}_{\chi}^{2}(\mathrm{G})\) 中）的一个子表示）。

我们来证明关于 \(\gamma_{\mathrm{G},\overline{\chi}}\) 的断言，第二个断言的证明是类似的。

 先设表示 \(\gamma_{\mathrm{G},\overline{\chi}}\) 存在一个同构于 \(\pi\) 的子表示 E。由于空间 E 非零，存在函数 \(f_1 \in \mathcal{K}_{\overline{\chi}}(\mathrm{G})\)，其类 \(\widetilde{f}_1\) 不与 E 正交。我们有 \(f_1 \in \mathcal{L}_{\overline{\chi}}^1 (\mathrm{G})\)。用 \(\widetilde{f}_{1,\mathrm{E}}\) 表示 \(\widetilde{f}_1\) 在 E 上的正交投影；它是 E 的一个非零元素。再设非零的 \(\widetilde{f}_2 \in \mathrm{E}\)。映射 \(h: g \mapsto \langle \widetilde{f}_{1,\mathrm{E}} | \gamma_{\mathrm{G},\overline{\chi}}(g)\widetilde{f}_2 \rangle\) 是 \(\pi\) 的一个矩阵系数。对每个 \(g \in \mathrm{G}\)，我们有 \(h(g) = \langle \widetilde{f}_1 | \gamma_{\mathrm{G},\overline{\chi}}(g)\widetilde{f}_2 \rangle\)，因为 \(\widetilde{f}_1 - \widetilde{f}_{1,\mathrm{E}}\) 与 E 正交。由 V 第 418 页的引理 8，函数 \(h\) 属于 \(\mathcal{L}_{\chi}^{2}(\mathrm{G})\)，于是命题 6 蕴含 \(\pi\) 模中心平方可积。

反之，设 \(\pi\) 模中心平方可积。设 \(x_0 \in \mathrm{E}\) 是一个非零向量。由命题 7 与公式 (10)，映射 \(y \mapsto \check{f}_{x_0,y}\) 是 \(\pi\) 到 \(\gamma_{\mathrm{G},\overline{\chi}}\) 中的一个单射 G-态射。

定义 4. --- 设 Z 是 C 的一个使 C/Z 为紧的闭子群。设 π 是 G 的一个模中心平方可积的不可约酉表示。满足命题 7 之性质的唯一实数 c \textgreater{} 0 称为 π 相对于 Z 的形式次数。把它记作 \(c_{\mathrm{Z}}(\pi)\)。

形式次数依赖于 Z 上哈尔测度的选取。若把 Z 上的哈尔测度 \(\beta_{Z}\) 乘以一个实数 t \textgreater{} 0，则 G/Z 上的测度 \(\nu_{Z} = \mu/\beta_{Z}\) 被乘以 \(t^{-1}\)，而 \(\pi\) 的形式次数被乘以 t。

形式次数由下述性质所刻画：

命题 8（正交关系）. --- 设 π 是 G 在希尔伯特空间 E 中的一个模中心平方可积的不可约酉表示。则有

\[
c _ {\mathrm{Z}} (\pi) \int_ {\mathrm{G} / \mathrm{Z}} \overline {{\langle x | \pi (g) x ^ {\prime} \rangle}} \langle y | \pi (g) y ^ {\prime} \rangle d \nu_ {\mathrm{Z}} (g) = \overline {{\langle x | y \rangle}} \langle x ^ {\prime} | y ^ {\prime} \rangle
\]

对所有 \((x,y,x',y')\in \mathrm{E}^4\) 成立

用 \(w\) 表示命题 7 中的态射。我们有

\[
\int_ {\mathrm{G} / \mathrm{Z}} \overline {{\langle x | \pi (g) x ^ {\prime} \rangle}} \langle y | \pi (g) y ^ {\prime} \rangle d \nu_ {\mathrm{Z}} (g) = \langle f _ {x, x ^ {\prime}} | f _ {y, y ^ {\prime}} \rangle
\]

而根据出处同上，由此得出

\[
\begin{array}{c} \langle f _ {x, x ^ {\prime}} | f _ {y, y ^ {\prime}} \rangle = \frac {1}{c _ {\mathrm{Z}} (\pi)} \langle w (x \otimes x ^ {\prime}) | w (y \otimes y ^ {\prime}) \rangle \\ = \frac {1}{c _ {\mathrm{Z}} (\pi)} \langle x \otimes x ^ {\prime} | y \otimes y ^ {\prime} \rangle = \frac {1}{c _ {\mathrm{Z}} (\pi)} \langle x | y \rangle_ {\overline {{\mathrm{E}}}} \langle x ^ {\prime} | y ^ {\prime} \rangle_ {\mathrm{E}}, \end{array}
\]

由此即得结论。

除上述命题外，对于互不同构的平方可积不可约表示，我们还有下述关系。

命题 9. --- 设 \(\pi_1\) 与 \(\pi_2\) 是 G 在希尔伯特空间 E \(_1\) 与 E \(_2\) 中的两个互不同构的不可约酉表示。假设 \(\pi_1\) 与 \(\pi_2\) 都模中心平方可积，且它们的中心特征标到 Z 的限制相同。则有

\[
\int_ {\mathrm{G} / \mathrm{Z}} \overline {{\langle x \mid \pi_ {1} (g) x ^ {\prime} \rangle}} \langle y \mid \pi_ {2} (g) y ^ {\prime} \rangle d \nu_ {\mathrm{Z}} (g) = 0
\]

对所有 \((x,x^{\prime},y,y^{\prime})\in \mathrm{E}_{1}^{2}\times \mathrm{E}_{2}^{2}\) 成立

 对 \(i=1,2\)，用 \(w_{i}\) 表示关于表示 \(\pi_{i}\) 的命题 7 中的态射。由 V 第 384 页的引理 8 \(b\) ) 与 V 第 394 页的命题 8 的断言 \(b\) )，\(w_{i}\) 的像含于 \(\pi_{i}\) -同型分量（\(\delta_{\mathrm{G},\chi}\) 的）中。由出处同上的断言 \(a\) )，\(w_{1}\) 的像因此与 \(w_{2}\) 的像正交。于是得出 \(\langle w_{1}(x\otimes x^{\prime})|w_{2}(y\otimes y^{\prime})\rangle=0\)，对所有 \((x,x^{\prime},y,y^{\prime})\in\mathrm{E}_{1}^{2}\times\mathrm{E}_{2}^{2}\) 成立，这就是所要证的公式。

注. --- 命题 8 与命题 9 的正交关系推广了 A, VIII, 第 399 页中的正交关系（另见 V 第 457 页 § 2 中 G 为紧的情形）。

